\section{Discussion}
\label{sec:discussion}

\paragraph{Fee spikes do not break channels}
A common view holds that fee spikes break Lightning channels: commitment fees signed in calm markets fall behind,
peers disagree over \texttt{update\_fee}, and channels are forced closed in a wave. In confirmed force closes we see no
such wave, for pre-anchor channels as little as for anchor channels. \NSpikesUp{} of the \NSpikes{} spikes did bring
more force closes, the most in May~2023, but pooled they sit below ordinary weeks of the same years.
Spikes still matter, but through what happens after a close: HTLC transactions confirm later, owners bump commitments,
and sweepers wait. The caveat is that the chain records confirmation, not broadcast. A close broadcast in a spike at
too low a fee may wait for a cheaper block, which would spread a wave of broadcasts over the following weeks; those
weeks show no rise beyond placebo noise either.

\paragraph{Where the anchor design worked, and where it did not}
Anchor outputs were meant to let the broadcaster pay the fee the market asks at broadcast time. For commitments that
is rarely needed: outside 2023, owners seldom bumped, and in 2026 most same-block owner children did not raise the fee
rate at all. Where the design shows its value is in HTLC transactions. In fee spikes, anchor channels' HTLC timeouts
were late far less often than pre-anchor ones, and late HTLCs are what costs a forwarding node money. Meanwhile the
anchors themselves became third-party dust: owners leave them, sweepers collect them at the first block or wait for
cheap blocks, and miners take most of their value. This is the behaviour LND's open issue on dangling anchors
assumes~\cite{lnd8997}, measured, but only for P2WSH anchors. No third party sweeps simple-taproot anchors, and their
owners rarely do: of the \TaprootAnchorsTwentySix{} created in 2026, \TaprootAnchorsSpentTwentySix{} were spent. Each
taproot force close therefore leaves outputs in the UTXO set, which a zero-value pay-to-anchor output, spent together
with its commitment, does not (\S\ref{sec:zerofee}).

\paragraph{The migration}
Zero-fee commitments change little in block space. Their gains are a simpler fee story, no 330-sat outputs to sweep,
and relay rules that no longer depend on the removed carve-out. Their test is the one on which anchors improved on
the old format: whether HTLC transactions confirm before their deadlines when fees spike. Our measurements are the
baseline for it: at 50~sat/vB or more, \HtlcLateSpikeAnchor\% of anchor channels' HTLC timeouts confirmed more than
\HtlcMarginTight{} blocks late. How fast these gains arrive depends on operators,
not on the specification: a channel keeps its format until it closes. Outside one cluster of short-lived channels,
pre-anchor commitments were \LegacyPctExClusterTwentySix\% of 2026's public force closes, in channels a median of
\LegacyOldAgeYears~years old, years after anchors became available. Our second paper found that most Lightning
operators upgrade within months of a release while a tail stays far behind~\cite{kumar2026upgradedynamics}, so the old
formats are likely to stay on-chain for years. The same choice, between fixing a fee when a transaction is signed and paying it
when it is broadcast, faces every protocol built on pre-signed transactions and timelocks, and there too what matters
is whether the time-critical transactions confirm on time~\cite{riard2020timedilation}.

\paragraph{What a force close costs}
In fees, a force close is cheap and a channel's life costs a fraction of a percent of its capacity, and the open is a
large part of that. The cost of a force close is a ratio, a multiple of a cooperative close in fees and block space,
and a wait, days before the closer's own funds return; the wait outweighs the fees only in low-fee years. Its
exposure, the HTLCs that must confirm before their deadlines while fees are high, has rarely become a loss: since
\SplitResFromYear, forwarding nodes lost \SplitResForwarderPayments{} payments that way, each with the preimage already
on-chain.

\paragraph{What the chain can catch}
Two 2026 events show what transactions alone can reveal. A week of revoked commitments matched a splicing bug that
Core Lightning had fixed days before, and the only split resolutions in the order replacement cycling leaves were a
test of that attack on mainnet. We found both without node logs, and the same tests can be run on new blocks as they
arrive.

\section{Limitations}
\label{sec:limitations}
\emph{Broadcaster.} The chain does not show which party broadcast a force close, so costs are per channel, both
parties together.
\emph{Invisible bumps.} Fee bumps that never confirmed, and replaced transactions, leave no trace on the chain; we
have no mempool record.
\emph{Private channels.} Private channels appear only when they force close, so their number and their force-close
rate are unknown, and their cooperative closes are not in the closing bill. Before 2024, ``unlisted'' also includes
public channels missing from our list.
\emph{Open channels.} Channels whose nodes have left count as open until closed on-chain, which lowers the
per-channel rate as they accumulate; cohort incidence is not affected.
\emph{Penalties.} Only punished breaches are visible.
\emph{Scope of the cost account.} Pre-anchor \texttt{to\_remote} spends are excluded by design, and third-party
anchor sweeps are counted as value lost, not as fees the channel pays. The break-even rate counts the broadcaster's
locked balance only, not pending HTLC amounts. Simple taproot channels' spends are not
classified, so their delays, HTLCs and bumps are missing from \S\ref{sec:cost}; they were \UnlTaprootPctTwentySix\% of
2026's unlisted force closes.
\emph{Fee regimes.} Every result depends on the fee regime, which is why the window covers all of Lightning's history.
Anchor channels appear in our data from 2021, so anchor results start then.
\emph{Lifecycle.} Opens are costed under an attribution rule; dual-funded opens are not identified (which does not
change a per-channel cost), and only completed lives enter, so short-lived channels weigh more.
\emph{Client labels} come from our first paper, inherit its accuracy and start in May~2023.
\emph{Association.} The fee-spike test is a within-year association; it rules out a large same-week effect, not every
causal path.

\section{Ethics}
\label{sec:ethics}
All inputs are public blocks and public gossip. We report aggregates only: no list, ranking or label of nodes, sweeping
operators, addresses or transactions, in the paper or in the released data. In particular we do not identify the
dominant anchor sweeper, the nodes of the first zero-fee commitment closes, the operator of the replacement-cycling
test or the channels of the breach wave.
The paper reports no new vulnerability: replacement cycling was disclosed in 2023~\cite{riard2023replacement}, and the
splicing bug was fixed before the wave~\cite{cln26068}.

\section{Conclusion}
\label{sec:conclusion}
Fee spikes do not bring waves of Lightning force closes, for channels whose fee is fixed at signing or for those that
can bump it.
What a force close costs is a multiple of a cooperative close in fees and days of locked funds, and the wait outweighs
the fees only in low-fee years. When fees are high, pending payments are exposed, yet forwarders lost few of them, each
with the preimage already on-chain. The anchor design reduced that exposure for HTLCs, while its P2WSH anchors became
dust that third parties sweep and miners collect, and its taproot anchors stay in the UTXO set. Zero-fee commitments
will save little block space; the measure of the migration is whether payments still resolve before their deadlines
when the next spike comes.
